\documentclass[10pt,a4paper]{article}

\usepackage{main}
\usepackage{exercices}

\renewcommand{\typedocument}{Exercices}
\renewcommand{\classe}{3e}
\renewcommand{\titre}{Révision Brevet}

\newcounter{numqglobal}

\newcommand{\q}[1]{%
    \refstepcounter{numqglobal}%
    \vspace{0.15em}%
    \noindent
    \thenumqglobal.\;#1%
    \par
}

\begin{document}

\creertitre{Révision Brevet --- Physique-Chimie}

\begin{multicols*}{2}

% ============================================================
%  CHIMIE
% ============================================================

\exercice{Constitution de la matière}

\q{Classer chacune des entités suivantes entre atome, ion ou molécule : CH\textsubscript{4}, H, dioxyde de carbone, Fe\textsuperscript{2+}, dioxygène, O\textsuperscript{2-}, H\textsubscript{2}O}

\q{Nommer les particules contenues dans le noyau d'un atome, et celles qui se déplacent autour du noyau.}

\q{L'atome de fer est noté $^{56}_{26}$Fe. Donner son nombre de protons, d'électrons et de neutrons.}

\q{Donner le nom et le nombre d'atomes présents dans la molécule d'octane C\textsubscript{8}H\textsubscript{18}.}

\q{Un atome de zinc Zn perd deux électrons. Écrire la formule de l'ion ainsi formé.}

\exercice{Transformations chimiques}

\q{On synthétise de l'ammoniac à partir de diazote et de dihydrogène :
\[ \text{N}_2 + 3\,\text{H}_2 \longrightarrow 2\,\text{NH}_3 \]
On dispose de \SI{824}{kg} de diazote et \SI{176}{kg} de dihydrogène. Quelle sera la masse d'ammoniac produite ?}

\q{Le pH d'une solution est égal à \num{9,3}. La solution est-elle acide ou basique ? Quels ions sont majoritaires (H\textsuperscript{+} ou HO\textsuperscript{-}) ?}

\q{Compléter les coefficients de la réaction de combustion du méthane en remplaçant chaque \textbf{?} par le bon nombre :
\[ \text{C}\text{H}_{4} + \mathbf{?}\,\text{O}_2 \longrightarrow \mathbf{?}\,\text{CO}_2 + \mathbf{?}\,\text{H}_2\text{O} \]}

\q{Rédiger le protocole expérimental d'un test caractéristique.}

\exercice{Propriétés de la matière}

\q{On chauffe un glaçon jusqu'à ce qu'il devienne de l'eau liquide. Comment appelle-t-on cette transformation physique ? La masse et le volume de l'eau ont-ils changé ?}

\q{Un bloc de métal a une masse de \SI{40,5}{g} pour un volume de \SI{15}{mL}. Calculer sa masse volumique en précisant l'unité. Pour identifier le métal, faut-il utiliser la masse ou la masse volumique ? Pourquoi ?}

% ============================================================
%  MÉCANIQUE
% ============================================================

\exercice{Mouvements et vitesse}

\q{Donner la définition d'un mouvement rectiligne accéléré et d'un mouvement circulaire uniforme.}

\q{Rappeler la formule de la vitesse moyenne d'un objet, en précisant les unités.}

\exercice{Interactions et forces}

\q{Une bûche de bois flotte sur l'eau. Quelles sont les forces qui s'exercent sur la bûche ? Dessiner le diagramme objet-interaction.}

\q{Quelle est la masse d'un rhinocéros dont le poids sur Terre est de \SI{20000}{N} ? (On prendra $g = \SI{10}{N/kg}$)}

\q{Représenter le rhinocéros par un rectangle et son poids par une flèche, en utilisant l'échelle : \mbox{1 cm $\leftrightarrow$ \SI{4000}{N}}.}

% ============================================================
%  ÉNERGIE
% ============================================================

\exercice{Formes et conversions d'énergie}

\q{Citer les 6 formes d'énergie vues en cours.}

\q{Donner la définition d'une source d'énergie renouvelable.}

\q{S'il n'y a aucun échange d'énergie avec l'extérieur, que peut-on dire de l'énergie totale d'un système ?}

\q{Lors d'une chute libre, en quelle forme d'énergie est convertie l'énergie potentielle de pesanteur ?}

\q{Si l'altitude d'un objet diminue, que se passe-t-il pour son énergie potentielle de pesanteur ? Rappeler la formule de $E_{pp}$ avec les unités de chaque grandeur.}

\q{Calculer l'énergie cinétique d'un rhinocéros de masse $m = \SI{3000}{kg}$ chargeant à $v = \SI{54}{km/h}$. Donner le résultat en joules.}

\q{Le rayonnement est-il un mode de transfert d'énergie ou de conversion d'énergie ?}

% ============================================================
%  ÉLECTRICITÉ
% ============================================================

\exercice{Puissance et énergie électrique}

\q{Quelle est la loi d'additivité des intensités ? Dans quel type de circuit s'applique-t-elle ?}

\q{Rappeler la loi d'Ohm et préciser les unités de chaque grandeur.}

Une lampe fonctionne avec une puissance de \SI{44}{W} sous une tension de \SI{220}{V}.

\q{Calculer l'intensité du courant qui la traverse.}

\q{Si la lampe est allumée pendant \SI{8}{h}, quelle énergie électrique a-t-elle consommée ? Préciser l'unité.}


% ============================================================
%  SIGNAUX SONORES
% ============================================================

\exercice{Signaux sonores et lumineux}

\q{Quelles sont les fréquences audibles pour l'être humain ?}

\q{Des dauphins utilisent un sonar (ondes sonores) pour détecter les obstacles. La vitesse du son dans l'eau est $v = \SI{1500}{m.s^{-1}}$. Combien de temps s'écoule-t-il entre l'émission et la réception d'une onde si l'obstacle est à \SI{500}{m} ?}

\q{Par quoi peut-on représenter un rayon lumineux ?}

\end{multicols*}
\end{document}